%% Unit 3 -- Plan before you build.  (Draws from Ch5 Planning and Alignment.)
%% Restyled to the sparse, one-idea-per-slide house style of units 1-2:
%% one thesis-shaped title, prose-first bodies, EXACTLY one red per rendered
%% frame (the two cmdtable frames keep their \gcmd slash-commands as identity).
%% Scientific claims are cited via \slidesources; the closing frame prints the
%% unit's papers. No preamble here; \input by the master (disciplined_vibe.tex).
%%
%% ORDER follows the software-engineering process:
%%   A. Motivation      -- why naive vibe coding breaks (four failure modes).
%%   B. Tooling setup   -- model card; agents/subagents; models+agents cmds;
%%                         plan vs build modes; the skill mechanism (AGENTS.md,
%%                         a skill is a file).
%%   C. Plan-arc map    -- the plan catalogue in four columns (align, model, scope + slice, spike).
%%   D. Plan arc, in process order -- askuserquestion (align),
%%                         domain-modeling (name), to-prd (scope),
%%                         to-issues (slice), prototype (spike a design question),
%%                         codebase-design (deep module).
%%   E. Breakout        -- plan the banana, write no code.
%% The build half (handoff + TDD, code-review, diagnosing-bugs,
%% improve-codebase-architecture, write-a-skill) now lives in Unit 4.

\sectionsubtitle{On Humansplaining and LLM Understanding}
% ======================================================================
\section[Alignment]{Plan}
% ======================================================================


% ======================================================================
% A. MOTIVATION
% ======================================================================

% === Vibe coding is brittle: the four failure modes ===
\begin{frame}{Pareto hits the vibe. Hard.}
  \slidesources{osmani2024seventy}

  % TODO (light Pocock spine): Pocock's catalogue exists to fix four failures:
  % misalignment (you and the agent mean different things), verbosity (400 lines
  % you never asked for), non-functional code, and architectural decay. Block 2
  % fixes each with a discipline. Unit 3 takes the first two; Unit 4 the last two.
  % Source: Pocock skills catalogue (pocock2025skills, four failure modes).
  \vskip0.3em
  {\small Naive vibe coding is brittle, then fails. Software engineering
  discipline (and a better model) fixes each mode.}
  \begin{columns}[T,onlytextwidth]
    \begin{column}{0.23\textwidth}
      \begin{tdblock}[equal height group=vibefail]
        \small
        \alertbf{Misalignment.}\; you and the agent mean different things by the
        same word.
      \end{tdblock}
    \end{column}
    \begin{column}{0.23\textwidth}
      \begin{tdblock}[equal height group=vibefail]
        \small
        \alertbf{Verbosity.}\; you ask for one thing and get 400 lines, while
        other functions silently break.
      \end{tdblock}
    \end{column}
    \begin{column}{0.23\textwidth}
      \begin{tdblock}[equal height group=vibefail]
        \small
        \alertbf{Non-functional.}\; the high-score save looks wired, but nothing
        persists.
      \end{tdblock}
    \end{column}
    \begin{column}{0.23\textwidth}
      \begin{tdblock}[equal height group=vibefail]
        \small
        \alertbf{Decay and Debt.}\; each build-loop makes project and module
        complexity harder to manage.
      \end{tdblock}
    \end{column}
  \end{columns}
  \end{frame}


% ======================================================================
% B. TOOLING SETUP
% ======================================================================

  \begin{frame}{Reading a model card}
    \begin{withmargin}
      \begin{tdmain}
        \centering
        \resizebox{\linewidth}{!}{%
        \begin{tikzpicture}[
            node distance=2mm,
            seg/.style={draw=tdfg, thick, rounded corners=2pt,
              inner xsep=7pt, inner ysep=6pt, text=tdfg, font=\ttfamily\large},
            seghi/.style={seg, draw=DHBWred, text=DHBWred},
            sep/.style={text=tddim, font=\ttfamily\large},
            lbl/.style={text=tddim, font=\scriptsize, align=center}]
          \node[seg] (fam) {qwen2.5};
          \node[sep, right=of fam] (c1) {-};
          \node[seg, right=of c1] (dom) {coder};
          \node[sep, right=of dom] (c2) {:};
          \node[seghi, right=of c2] (sz) {7b};
          \node[sep, right=of sz] (c3) {-};
          \node[seg, right=of c3] (pt) {instruct};
          \node[sep, right=of pt] (c4) {-};
          \node[seg, right=of c4] (q) {q4\_K\_M};
          \node[lbl, below=2.5mm of fam] {family\\+ version};
          \node[lbl, below=2.5mm of dom] {domain};
          \node[lbl, text=DHBWred, below=2.5mm of sz]  {parameters};
          \node[lbl, below=2.5mm of pt]  {post-training};
          \node[lbl, below=2.5mm of q] {quantization};
        \end{tikzpicture}}
        \vskip1.1em
        \begin{tdblock}
          {\small Two things the name never carries. The \textbf{context window}
          (\texttt{num\_ctx}) is the token budget the model attends to at once,
          paid for in \textbf{GPU memory}. And whether the model \textbf{reasons}
          on a scratchpad before it answers.}
        \end{tdblock}
        \vskip0.5em
        {\footnotesize\color{tddim}Feel free to switch to more capable models now.}
      \end{tdmain}
      \begin{tdmargin}
        {\color{tdfg}\textbf{qwen2.5}}\; base model and release.\par\smallskip
        {\color{tdfg}\textbf{coder}}\; fine-tuned for a domain, here code.\par\smallskip
        {\color{DHBWred}\textbf{7b}}\; \alert{parameters}: seven billion weights
        in memory (15GB) while the model runs.\par\smallskip
        {\color{tdfg}\textbf{instruct}}\; follows instructions and chats; a
        \texttt{base} model only continues text.\par\smallskip
        {\color{tdfg}\textbf{q4\_K\_M}}\; quantization: about 4 bits per weight,
        shrinking memory to 4.4\,GB.
      \end{tdmargin}
    \end{withmargin}
  \end{frame}

% === Agents and subagents (concept, pairs with the cmdtable below) ===
\begin{frame}{An agent is a prompt, its tools, and a model}
  \begin{withmargin}
    \begin{tdmain}
      \textbf{Plan} and \textbf{build} are the two built-in \textbf{agents}. Each
      bundles a system prompt, a set of allowed tools, and a model. You define
      your own the same way: one markdown file in \texttt{.opencode/agent/}.
      \vskip0.8em
      A \alertbf{subagent} is an agent a primary one spawns to handle a single
      self-contained task in its \textbf{own fresh context window}, then hands
      back a short summary. The side quest never clutters your main session.
      \vskip0.9em
      \begin{tdblock}
        {\ttfamily\small
        > @007 shake it, do not stirr.}
      \end{tdblock}
    \end{tdmain}
    \begin{tdmargin}
      Switch primary agents with \texttt{/agent} or \texttt{Tab}; call a subagent
      by name with \texttt{@}.
      \par\medskip
      Delegation is a tool call. Enter agent orchestration.
    \end{tdmargin}
  \end{withmargin}
\end{frame}

% === Models and agents (cmdtable) ===
\begin{frame}{Models and agents}
  \begin{cmdtable}
    \gcmd{/models}\gsep\gkey{ctrl+x m} &
      Switches the model the agent runs on, from your configured providers.\\[4pt]
    \gkey{f2} &
      \textbf{Recent model:} toggles back to the model you used just before, for
      a quick side-by-side on the same task.\\[4pt]
    \gcmd{/agent}\gsep\gkey{ctrl+x a} &
      Switches the active agent, build versus plan, each with its own tools and
      system prompt.\\[4pt]
    \gkey{tab} &
      \textbf{Cycle agent:} steps to the next agent without opening the list.\\[4pt]
    \gkey{right} &
      \textbf{Cycle subagents:} moves between the subagents a session has spawned,
      to read what each is doing.\\
  \end{cmdtable}
\end{frame}

% -------------------------------------------- plan vs build (two modes)
\begin{frame}{Plan and Build mode in opencode}
  \vskip0.3em
  {\small Plan mode reads and reasons; Build mode writes and runs. Like gating,
  the split mitigates surprises. Toggle with \texttt{Tab}.}
  \vskip0.8em
  \begin{columns}[T,onlytextwidth]
    \begin{column}{0.47\textwidth}
      \begin{tdblock}
        \small
        \textbf{Plan}\par\smallskip
        \alertbf{Read only}. opencode explores, explains, and proposes; it never
        edits files or runs commands.
        \vskip0.5em
        Use it to scope the work and agree on the approach before any change lands.
      \end{tdblock}
    \end{column}
    \begin{column}{0.47\textwidth}
      \begin{tdblock}
        \small
        \textbf{Build}\par\smallskip
        \alertbf{Full access}. opencode edits files and runs tools to carry out
        a plan. Here small models default into action.
        \vskip0.5em
        Use it once the plan is clear to both parties. You can build across
        several sessions.
      \end{tdblock}
    \end{column}
  \end{columns}
  \vskip0.9em
  \centering
\end{frame}


% === A skill is just a file the agent reads ===
\begin{frame}{Agentspanning-Capabilities or Skills}
  \begin{withmargin}
    \begin{tdmain}
      A \alertbf{skill} is a markdown file of instructions and context. Invoke it
      by name; the agent opens that file and follows it, turn by turn. That is the
      whole mechanism.
      \vskip0.6em
      Those instructions do more than steer the model. A skill can call an API,
      run a script, or hand a slice of work to a subagent, so one named file
      becomes the interface to agentic workflow automation.
    \end{tdmain}
    \begin{tdmargin}
      Invoke with /skills, then pick one.
      \par\medskip
      Compact, token-frugal skills win: the whole file is read into context every
      time the skill runs.
    \end{tdmargin}
  \end{withmargin}
\end{frame}

% === write-a-skill ===
\begin{frame}{\texttt{write-a-skill}}
  \begin{withmargin}
    \begin{tdmain}
      Every skill so far taught the agent one job. \texttt{write-a-skill} teaches
      it to author the next one in the same shape, so the toolkit grows by its own
      rules. Name one job, list \textbf{3 to 5 bold-lead steps}, close on a single
      principle, all short enough to read in a breath. If the job needs an ``and'', split it,
       and wrap it in another skill.
      \alertbf{skillception}.
      \vskip0.6em
      \begin{tdblock}
        {\ttfamily\small
        > write-a-skill for "writing skills" \\
        Or wait, I will just nano this myself.}
      \end{tdblock}
    \end{tdmain}
    \begin{tdmargin}
      Writing a skill, is a skill. Of course it is easy to forget, that we can also still do stuff ourselves - manually.
    \end{tdmargin}
  \end{withmargin}
\end{frame}


% ======================================================================
% C. PLAN-ARC MAP
% ======================================================================

% === the plan toolkit: the read-only catalogue this unit walks, one frame ===
% Only the plan arc lives here. The build skills (tdd, code-review,
% diagnosing-bugs, improve-codebase-architecture) and the greenfield on-ramp
% belong to Unit 4; keep them off this frame so it maps only what Unit 3 covers.
\begin{frame}{The plan toolkit, read only}
  \vskip0.3em
  \begin{columns}[T,onlytextwidth]
    \begin{column}{0.23\textwidth}
      \begin{tdblock}
        \scriptsize
        \alertbf{Align}\par\smallskip
        \texttt{askuserquestion}
        \par\medskip
        Reach a design you both hold before a line is written.
      \end{tdblock}
    \end{column}
    \begin{column}{0.23\textwidth}
      \begin{tdblock}
        \scriptsize
        \alertbf{Model}\par\smallskip
        \texttt{domain-modeling}
        \par\medskip
        Fix one word to one meaning, then name the pieces.
      \end{tdblock}
    \end{column}
    \begin{column}{0.23\textwidth}
      \begin{tdblock}
        \scriptsize
        \alertbf{Scope + slice}\par\smallskip
        \texttt{to-prd}\par
        \texttt{to-issues}
        \par\medskip
        Write the brief, then split it into small issues.
      \end{tdblock}
    \end{column}
    \begin{column}{0.23\textwidth}
      \begin{tdblock}
        \scriptsize
        \alertbf{Spike}\par\smallskip
        \texttt{prototype}
        \par\medskip
        Throw one away to settle a single design question.
      \end{tdblock}
    \end{column}
  \end{columns}
\end{frame}


% ======================================================================
% D. PLAN ARC, IN PROCESS ORDER
%    align -> name -> scope -> slice -> spike -> design
% ======================================================================

% === askuserquestion: Align before you name ===
\begin{frame}{\texttt{askuserquestion}}
  \begin{withmargin}
    \begin{tdmain}
      The arc starts with \alertbf{askuserquestion}. It does not write a plan; it
      reaches one. The skill interviews you one question at a time, proposes its
      own answer to each and says why, and walks every branch of the decision tree, until it has
      no questions left.
      \vskip0.6em
      \begin{tdblock}
        {\ttfamily\small
        > run askuserquestion on feature\_brief.md\\
        Q: do power-ups stack? (I'd say no)\\
        > no\\
        Q: does score reset on death?\,\ldots\\
        no open questions. aligned.}
      \end{tdblock}
    \end{tdmain}
    \begin{tdmargin}
      One question at a time, and it stops only when alignment is reached, not
      before. In general following proposals is adviced, until it is not. 
      \par\medskip
      And just because a model has no questions left, does not mean everything is perfectly clarified.
    \end{tdmargin}
  \end{withmargin}
\end{frame}

% === domain-modeling: One word, one meaning, everywhere ===
\begin{frame}{\texttt{domain-modeling}}
  \slidesources{evans2003ddd}
  \begin{withmargin}
    \begin{tdmain}
      \texttt{askuserquestion} settled the vocabulary. You have no
      \texttt{GLOSSARY.md} yet; \texttt{domain-modeling} rereads the chat and code
      and writes one line per term. Now code and tests speak the same words every
      session. \textcolor{DHBWred}{Naming the thing is the alignment.}
      \vskip0.6em
      \begin{tdblock}
        {\ttfamily\small
        > domain-modeling}
      \end{tdblock}
    \end{tdmain}
    \begin{tdmargin}
      \begin{tdcodet}{GLOSSARY.md (after the run)}
        obstacle: box to clear\\
        hitbox: what collides()\\
        duck: lowers the hitbox\\
        ramp: speeds up w/ score
      \end{tdcodet}
      \par\smallskip
      Remember how the substrate of these models is similarity within an embedded feature space.
      So whenever possible go with their wording to increase performance and save context tokens.
    \end{tdmargin}
  \end{withmargin}
\end{frame}

% === to-prd: Write down what you're not building ===
\begin{frame}{\texttt{to-prd}, the product requirements document}
  \begin{withmargin}
    \begin{tdmain}
      Start with letting the agent write this document based on the survey, then step in and adjust, extend, and clean up.
      \vskip0.8em
      \begin{tdblock}
        {\ttfamily\small
        Acceptance lines\\
        - the high score survives a page reload\\
        - the score ticks up each step\\
        Out of scope\\
        - no multiplayer, no leaderboard\\
        - no power-ups this pass}
      \end{tdblock}
    \end{tdmain}
    \begin{tdmargin}
      to-prd fixes verbosity: the aligned chat becomes a short PRD in
      \texttt{issues/PRD.md}. Two parts carry the weight, the out-of-scope list and
      one \textcolor{DHBWred}{acceptance line} per slice, a checkable sentence like
      ``the high score survives a page reload''.
    \end{tdmargin}
  \end{withmargin}
\end{frame}

% === to-issues: Slices you can grab ===
\begin{frame}{\texttt{to-issues}}
  \begin{withmargin}
    \begin{tdmain}
      \texttt{to-issues} splits \texttt{issues/PRD.md} into independently
      \alertbf{grabbable} slices, one markdown file each. Every issue carries its
      acceptance line as its ``how tested''.
      \vskip0.9em
      \begin{tdblock}
        {\ttfamily\small
        issues/01-banana-on-screen.md\\
        \phantom{ }tested: the banana runs across the screen\\[3pt]
        issues/02-persist-high-score.md\\
        \phantom{ }tested: the high score survives a reload\\
        \phantom{ }blocked-by: 01}
      \end{tdblock}
    \end{tdmain}
    \begin{tdmargin}
      Reject a horizontal first slice and re-slice until the first one shows a
      number. Check it actually writes the slice.mds into \texttt{issues/}.
    \end{tdmargin}
  \end{withmargin}
\end{frame}

% === prototype: Spike it, learn it, delete it ===
\begin{frame}{Spike by \texttt{prototype}}
  \begin{withmargin}
    \begin{tdmain}
      What if you can not answer a design question? \texttt{prototype} answers it with a
      \alertbf{throwaway spike}. The answer is the deliverable; the scratch code is
      deleted. 
      \vskip0.9em
      {\small Will this library do the job, is the API fast enough, does the real
      data fit the schema? A 20-line script answers all three. Record ``library
      works, 120\,ms median, dates are strings'', delete the script.}
    \end{tdmain}
    \begin{tdmargin}
      The delete step is mandatory. A spike you keep and build on rots into decay:
      it was never designed to last. It was a standalone accepted cost.
    \end{tdmargin}
  \end{withmargin}
\end{frame}


\iffalse
% === codebase-design ===
\begin{frame}{\texttt{codebase-design}}
  \begin{withmargin}
    \begin{tdmain}
      The outcome is a small set of function signatures,
      names and arguments only, with the bodies left empty. Keep the logic
      separate from anything it draws or displays. Name the one job the module
      does, borrow the shared names from \texttt{GLOSSARY.md}, then write only the
      functions that expose that job and hide the work behind them.
      \vskip0.6em
      \begin{tdblock}
        {\ttfamily\small
        > design the core interface\\
        create(config)\\
        step(state, input)\\
        done(state)\\
        node --test  ok, no display}
      \end{tdblock}
      \vskip0.5em
      {\small A few pure functions carry the whole job. Each takes plain
      arguments and returns plain values; the display layer lives outside and
      calls in.}
    \end{tdmain}
    \begin{tdmargin}
      The test of depth: \alertbf{if the core cannot run under \texttt{node
      --test} with no display, the interface is too wide}. A wide draft leaks
      rendering into the core, thin functions that each need a screen to run.
      Catch it with the headless test, then shrink the interface until only the
      pure functions remain.
    \end{tdmargin}
  \end{withmargin}
\end{frame}
\fi

% ======================================================================
% E. BREAKOUT -- plan the banana, write no code.
% ======================================================================

% === Vibe an infinitely running banana ===
\begin{frame}{\alertbf{Vibe} an infinitely running banana}
  \sideimagecover[0.33]{banana_1}
  \sidecaption{The banana runner, planned this time.}
  \begin{sidebody}
    Same endless runner, now under this discipline. Run the whole planning chain
    on \texttt{feature\_brief.md}, back to back, and write no code.
    \begin{itemize}
      \item \texttt{askuserquestion} until it stops asking.
      \item \texttt{domain-modeling} into \texttt{GLOSSARY.md}.
      \item \texttt{to-prd}: Set up a product requirements document.
      \item \texttt{to-issues}: A first vertical slice that serves as a silver bullet.
    \end{itemize}
    \vskip0.4em
    The deliverable is a \alertbf{structured executable plan}. If the agent starts writing
    code, stop it.
  \end{sidebody}
\end{frame}


% References are collected in one shared package at the end of the deck
% (see disciplined_vibe.tex), not per unit.
